\section{Tensor products over the shared parameter algebra}
\label{app:tensor}

We give a more explicit two-factor proof of \cref{thm:quotient-tensor}.  Let
\(B_1=A[X]\), \(B_2=A[Y]\) for disjoint finite variable sets and let
\(I_1\lhd B_1\), \(I_2\lhd B_2\).  The canonical isomorphism
\begin{equation}
  \alpha:B_1\otimes_A B_2\xrightarrow{\sim}A[X,Y]
\end{equation}
sends \(f(X)\otimes g(Y)\) to \(f(X)g(Y)\).  Let \(J\) be the ideal of the
tensor product generated by the images of \(I_1\otimes B_2\) and
\(B_1\otimes I_2\).  Right exactness of tensor product gives canonical
surjections whose composite induces
\begin{equation}
 (B_1/I_1)\otimes_A(B_2/I_2)
 \xrightarrow{\sim}(B_1\otimes_AB_2)/J.
\end{equation}
Under \(\alpha\), the ideal \(J\) maps to
\(I_1A[X,Y]+I_2A[X,Y]\).  Therefore
\begin{equation}
 \frac{A[X,Y]}{I_1A[X,Y]+I_2A[X,Y]}
 \cong (B_1/I_1)\otimes_A(B_2/I_2).
\end{equation}

Associativity of tensor products supplies the finite iterated statement.  No
ordering of samples is mathematically preferred; different parenthesizations
are related by the canonical associators.

\subsection{Why base choice matters}

If one instead forms a tensor product over \(k\), then the two copies of
\(A=k[\theta]\) remain independent.  For example,
\begin{equation}
  k[\theta]\otimes_k k[\theta]
  \cong k[\theta^{(1)},\theta^{(2)}],
\end{equation}
whereas
\begin{equation}
  k[\theta]\otimes_{k[\theta]}k[\theta]
  \cong k[\theta].
\end{equation}
The latter is the correct model for a single shared trainable parameter vector.

\subsection{Real points do not commute with arbitrary projection claims}

The tensor identity is an identity of algebras.  It implies that the global
equations contain no hidden cross-sample polynomial except through the shared
base.  It does not imply
\begin{equation}
  \Proj_\theta V(I)=V(I\cap A)
\end{equation}
as equality of real sets without closure and real-algebraic qualifications.
This is why the semantic projection in \cref{prop:semantic-projection} is kept
separate from elimination ideals.
